\section{실험 결과 및 성능 분석}
\label{sec:실험결과}

본 장에서는 Lite-GLOBE 시뮬레이션을 통해 획득한 \method\ 계열 모델의 종합적인 성능 평가 결과를 분석한다.

\subsection{종합 성능 평가}
총 14개 평가용 시나리오에 걸쳐 측정된 평균 성능 메트릭의 비교 요약은 표~\ref{tab:overall_results_kci}에 기술되어 있다. \phaseTwelve\ 모델(본 제안 기법의 이전 단계 완성 모형)은 높은 전송 성공도와 짧은 단일 지연 시간을 균형 있게 달성함을 보여준다.

\begin{table}[htbp]
\centering
\caption{통합 라우팅 메트릭 비교 요약 (Phase 12 풀 런 및 외부 기준선)}
\label{tab:overall_results_kci}
\begin{tabular}{lcccccc}
\toprule
알고리즘군 & PDR (\%) & 데드라인 준수율 (\%) & 평균 지연 (s) & 지연 95\% (s) & 제어 부하 (Byte) & 에너지 소비 \\
\midrule
GPSR & 53.9 & 49.2 & 1.214 & 6.814 & 16 & 0.421 \\
Predictive Geo & 84.9 & 81.4 & 0.864 & 4.452 & 192 & 0.284 \\
Evo-QGeo & 84.3 & 80.6 & 0.892 & 4.521 & 256 & 0.292 \\
IQMR & 81.2 & 76.5 & 0.942 & 4.814 & 160 & 0.312 \\
DRAMA & 84.8 & 81.1 & 0.884 & 4.612 & 320 & 0.279 \\
\textbf{\phaseTwelve\ (Ours)} & \textbf{86.4} & \textbf{83.1} & \textbf{0.814} & \textbf{4.264} & \textbf{128} & \textbf{0.252} \\
\bottomrule
\end{tabular}
\end{table}

\subsection{전송 신뢰성 및 시간 엄수 성능}
표~\ref{tab:overall_results_kci}에 제시된 바와 같이, \phaseTwelve\ 모델은 가장 높은 86.4\%의 PDR을 확보하여 기존 지리적 라우팅 프로토콜인 GPSR 대비 32.5\%라는 큰 성공률 차이를 벌렸다. 목적지 거리 위주로만 동작하여 우회 경로에 실패하는 GPSR과 달리, 오프라인 교사의 우회 거동을 1-hop 피처 기반으로 정밀하게 묘사하고 있기 때문이다. 
아울러, 데드라인 준수율 지표 역시 PDR 결과와 동일한 증가 경향을 유지하고 있다. 이는 제안 기법이 홉 수나 큐 지연을 능동적으로 억제함으로써 버퍼 내에 대기하는 노드를 최소화하고 실시간 전송 한계(Deadline)를 강력히 보증하고 있음을 지시한다.

\subsection{단일 패킷 지연 및 95\% 꼬리 지연}
\phaseTwelve\ 모델의 평균 종단간 전송 지연 시간은 0.814초이며, 95\% 꼬리(tail) 지연은 4.264초로 나타났다. 이는 예측 지리(4.452초), Evo-QGeo(4.521초), DRAMA(4.612초)와 같은 고성능 기법들과 비교해서도 가장 지연 변동 폭이 적다. 
비교군 중 아블레이션 모델인 Phase 8(지리 잔차 KD 단독)은 지연 시간이 0.792초로 미세하게 짧게 측정되었다. 그러나 이는 Phase 8 모델이 링크 단절 임계 상태에 도달한 패킷들을 미리 판단하지 못하고 빠르게 누출/드롭하기 때문에 생존한 극히 짧은 패킷들의 수치만 반영된 통계적 왜곡 현상이다. 따라서 신뢰도(PDR)와 지연 성능을 연계해서 보았을 때, 제안 모델의 안정성이 가장 높다.

\subsection{실행 정보 풋프린트 오버헤드}
차홉 결정을 내릴 때 개별 노드가 실시간으로 요구하는 입력 및 전송 데이터 크기를 비교 분석하였다. GPSR은 오직 위치 벡터만을 처리하므로 가장 낮다. 예측 라우팅 기법들과 MARL 계열(DRAMA)은 주변과의 무선 토폴로지 통신 데이터 교환 때문에 각각 192, 256, 320 Byte 수준의 비교적 큰 자원을 소모한다. 반면 제안된 \phaseTwelve\ 모델은 교사 지식을 local MLP로 완벽히 이전하고, switch 기반 위험 판단 시에만 1-hop predictive feature를 선별 조회하므로 128 Byte의 적은 크기로 구동될 수 있다.

\subsection{시나리오별 상세 동향}
공간 장벽이 배치된 라우팅 홀 시나리오에서 GPSR은 PDR 50\% 미만으로 붕괴된 반면, 제안 기법은 1-hop onward stability 예측 피처를 통해 홀 경계를 사전에 감지하고 오프라인 교사의 우회 경로 기법을 성공적으로 모사하여 85\% 이상의 성공률을 유지하였다.
이동 방향이 교차하여 링크가 파괴되는 예측 단절 시나리오에서도, 위험 스위치 게이트 모듈이 링크 잔여 시간과 여유 큐를 바탕으로 즉시 대안 노드를 활성화해 패킷 누출을 최소화하였다.

\subsection{Phase 13 P+ 제안 기능의 구현 상태}
새롭게 고안된 Phase 13 \method\ 프레임워크의 채널 손실 대처 필터 및 에너지 가중치 logit 보정(Algorithm 1)은 기본 동작 훈육을 통과하였다. 현재 데이터셋 전체 시드에 대한 최종 리포팅 런이 진행 중이다.
